% ─────────────────────────────────────────────────────────────────────────────
% Section II: Background and Related Work
% Structure: brief's argument flow (naive RAG fails -> retrieval not enough ->
% post-hoc can't prevent -> SELF-RAG contrast -> the gap).
% Every statistic here is verified against the cited source. Only ONE numeric
% claim survived verification (+17.9% GSM8K, Wang et al. 2023). All other prior
% percentages could not be confirmed in-source and were made qualitative.
% ─────────────────────────────────────────────────────────────────────────────
\section{Background and Related Work}
\label{sec:related}

Trust Before Text sits at the intersection of retrieval-augmented
generation, hallucination mitigation, and abstention. This section
organizes the prior work around a single question: \textit{where, if
anywhere, does a RAG system verify its evidence before committing to an
answer?} We argue that existing systems verify either not at all, only on
the retrieval side, or only after generation, and that none installs a
hard, pre-generation gate.

\subsection{Standard RAG and Its Failure Modes}
Lewis et al.~\cite{lewis2020rag} established the retrieve-then-read
paradigm, in which a dense retriever supplies passages to a
sequence-to-sequence generator. This improved factual grounding over
purely parametric models but exposed a persistent tension: the generator
can bypass the retrieved evidence and fall back on its pre-training
knowledge. The survey of Gao et al.~\cite{gao2023rag_survey} catalogs the
resulting failure modes of naive RAG, including hallucination under
retrieved context, brittleness on noisy or contradictory documents, and
the absence of principled abstention.

These failures are not merely theoretical. Shuster et
al.~\cite{shuster2021retrieval} show that retrieval augmentation reduces
but does not remove hallucination in knowledge-grounded generation, and
the RAGTruth corpus of Wu et al.~\cite{wu2024ragtruth} demonstrates that
hallucination occurs even when the retrieved passages contain the correct
answer, confirming that retrieval success does not guarantee faithful
generation. Siriwardhana et al.~\cite{siriwardhana2023rag} find that
domain-specific fine-tuning reduces, but does not eliminate, the model's
tendency to prefer parametric knowledge over retrieved evidence,
particularly under vocabulary shift, while Rakin~\cite{rakin2024rag}
reports that end-to-end fine-tuned document-QA systems retain a residual
hallucination rate. Crucially, when documents disagree, Wolk et
al.~\cite{wolk2025clinical} show that Fusion-in-Decoder architectures
synthesize a false consensus from contradictory passages without ever
surfacing the contradiction, precisely the behavior a high-stakes
deployment cannot tolerate.

\subsection{Retrieval Improvements Are Not Enough}
A large body of work improves the retrieval stage through query expansion,
hybrid sparse-dense fusion, and cross-encoder reranking, and these
techniques raise average-case precision and recall on clean
benchmarks~\cite{yu2024evaluation}. However, they operate entirely on the
retrieval side and leave the generator unconstrained. A retrieval score
measures query-document similarity, not evidential answerability: a
high-similarity chunk need not contain the answer, and need not be
consistent with the other chunks retrieved alongside it. Fadeeva et
al.~\cite{fadeeva2023reeval} make this fragility explicit, showing that
retrieval-augmented models fail to isolate the relevant evidence and lack
any mechanism to abstain when the retrieved context is adversarially
perturbed or misleading. Approaches that add a lightweight
retrieval-quality evaluator to trigger corrective actions, such as
Corrective RAG~\cite{yan2024crag}, improve robustness to low-quality
retrieval, but they inherit the evaluator's own error rate and still judge
relevance probabilistically rather than verifying evidential sufficiency.
Better retrieval, in short, produces better candidate evidence, but it
does not decide whether that evidence is trustworthy enough to answer
from.

\subsection{Post-Hoc Detection Cannot Prevent Hallucination}
A complementary line of work detects hallucination \textit{after} the
model has generated. SelfCheckGPT~\cite{manakul2023selfcheck} samples
multiple generations and flags inconsistency as a black-box hallucination
signal; semantic-entropy methods~\cite{farquhar2024semantic} estimate
uncertainty over meaning-equivalent generations to the same end; and
Chain-of-Verification~\cite{dhuliawala2023cov} has the model draft
verification questions to check its own claims. A related family improves
answer stability by sampling several reasoning paths and voting over them:
self-consistency, for instance, raises chain-of-thought accuracy by a wide
margin on arithmetic reasoning, with a reported gain of $+17.9\%$ on
GSM8K~\cite{wang2023selfconsistency}. Yet all of these methods act on the
model's own probabilistic behavior. They can improve or flag reliability,
but they neither prevent unsupported content from being produced nor
guarantee that an answer is grounded in the retrieved evidence: the
unsupported content has already been generated by the time it can be
detected. Evaluation frameworks such as RAGAS~\cite{es2023ragas} likewise
measure faithfulness and context relevance at the level of end-to-end
behavior; they can report \textit{that} a pipeline hallucinated, but not
\textit{at which stage} it became unsafe, and they impose no gate that
would have stopped it.

\subsection{SELF-RAG and Its Limitations}
SELF-RAG~\cite{asai2023selfrag} is the closest prior work to our approach
and the most important point of contrast. It trains the model to emit
reflection tokens during generation that signal retrieval utility and
factual support, effectively interleaving verification with synthesis in a
single pass. The verification signal, however, is a learned soft behavior:
the reflection tokens are a probabilistic function of training, so under
distribution shift the model may emit positive utility tokens while
producing unsupported text. The safety property is a tendency of the
trained model, not a guarantee of the architecture. Trust Before Text
takes the opposite stance. Verification is completed by a deterministic
pipeline \textit{before} the language model is invoked, so a failed
verification means the model is never called and the unsafe answer is
never generated. The distinction is between a signal the model is trained
to produce and a gate the system is built to enforce.

\subsection{The Gap: No Pre-Generation Verification Gate}
Taken together, the literature verifies evidence in three places, none of
them sufficient. Retrieval-side methods assess similarity, not
answerability or consistency. Post-hoc detectors and end-to-end evaluators
act only after generation, when the unsupported content already exists.
Learned in-generation signals such as SELF-RAG's reflection tokens remain
probabilistic and can fail silently. Attribution studies underline the
consequence: Rashkin et al.~\cite{rashkin2023attribution} show that
post-hoc attribution frequently misaligns claims with their cited sources,
mixing supported and unsupported information without an intrinsic,
verifiable mapping. What no existing system provides is a hard
architectural boundary that blocks generation until the evidence has been
shown to be sufficient and non-contradictory. That pre-generation
verification gate is the gap Trust Before Text fills.
